% Generated draft: requires mathematical and source review.
\documentclass[11pt]{article}
\usepackage[margin=25mm]{geometry}
\usepackage{fontspec}
\setmainfont{Latin Modern Roman}
\usepackage{amsmath,amssymb,mathtools,unicode-math}
\setmathfont{Latin Modern Math}
\usepackage{enumitem,needspace,xurl,hyperref}
\setlength{\parindent}{0pt}
\setlength{\parskip}{5pt}
\setlength{\emergencystretch}{3em}
\newcommand{\N}{\mathbb{N}}
\newcommand{\Z}{\mathbb{Z}}
\newcommand{\Q}{\mathbb{Q}}
\newcommand{\R}{\mathbb{R}}
\newcommand{\C}{\mathbb{C}}
\newcommand{\F}{\mathbb{F}}
\newcommand{\A}{\mathbb{A}}
\newcommand{\PP}{\mathbb{P}}
\newcommand{\E}{\mathbb{E}}
\newcommand{\eps}{\varepsilon}
\newcommand{\dd}{\,\mathrm{d}}
\begin{document}
\section*{UnsolvedMath ID: 30000671}
\textbf{Source problem code:} catalogue entry
\subsection{Definitions and mathematical statement}
Let $(A,\mathfrak m)$ and $(B,\mathfrak n)$ be commutative Noetherian local rings with identity, complete and separated for their maximal-ideal-adic topologies. Thus the canonical maps
\[
A\longrightarrow\varprojlim_{r\geq1}A/\mathfrak m^r,
\qquad
B\longrightarrow\varprojlim_{r\geq1}B/\mathfrak n^r
\]
are isomorphisms. The question is whether
\[
\left[\forall r\geq1\ \exists\text{ a ring isomorphism }
\phi_r:A/\mathfrak m^r\longrightarrow B/\mathfrak n^r\right]
\quad\Longrightarrow\quad A\cong B.
\]
The isomorphisms $\phi_r$ are not assumed compatible with the quotient maps. A compatible system, satisfying
\[
\pi^B_{r+1,r}\circ\phi_{r+1}
=\phi_r\circ\pi^A_{r+1,r},
\]
would induce an isomorphism on inverse limits, but the hypothesis gives only a separate isomorphism at each order. No common coefficient field, prescribed residue-field identification or coefficient-preserving structure is imposed. Although these quotients are Artinian, they are not assumed to be finite sets. An isomorphism of the local rings automatically preserves their maximal ideals and all their powers, hence their adic topologies.
\subsection{Short English statement}
Are complete Noetherian local rings isomorphic whenever their truncations are isomorphic at every order? The truncation isomorphisms may vary independently with the order; their compatibility is the essential missing condition to an immediate inverse-limit argument.
\subsection{Sources}
\begin{itemize}
\item Supplied report: \emph{Model Theory and Groups} (2007), Oberwolfach Reports, DOI 10.4171/owr/2007/02, pp. 83--138; proposed year 2007. \url{https://doi.org/10.4171/owr/2007/02}.
\item Local UnsolvedMath registry, permanent ID 30000671, catalogue entry; created 25 August 2026, updated 16 September 2026. The corrected\_verified note reports comparison with the official report. The supplied 21 August 2026 literature triage labels the extracted question open and reports no located authoritative complete resolution. No additional external reference is supplied; this dated assessment is not independently verified here.
\end{itemize}
\end{document}
